% page 316 du fac-simile (image p-315 du scan)
% bloc 154.9 x 242.0 mm ; marge gauche 8.0 mm
\pgc{-12.700mm}{0.000mm}{
\l{\cel{5}308}
\l{}
\l{}
\l{\sou{kornico}. (\sou{ariktekt.}) Ornamento salianta qua kronizas la cir-}
\l{\cel{3}klo di la entablamento di edifico. Ornamento analoga, ek}
\l{\cel{3}petro, ligno, e c., qua existas cirkum plafono, chambro,}
\l{\cel{3}qua kronizas moblo, pordo., ec. - EFIRS.}
\l{}
\l{\sou{korniko}. (\sou{zool.})\cel{2}Speco di korvo, mikra. -L.\cel{1}\sou{corvus cornix}.}
\l{\cel{3}- FIS.}
\l{}
\l{\sou{korodar}. (\sou{trans.}) Rodar lente, per efikivo kemiala (maxim-}
\l{\cel{3}multa-kaze). - DEFIS.}
\l{}
\l{\sou{koroido.}\cel{2}(\sou{anat.)}\cel{2}Membrano interna, kovrita ye materio ni-}
\l{\cel{3}gra (pigmento) qua tapisas la okulo. -}
\l{}
\l{\sou{korolo.}\cel{2}(\sou{bot.}) Parto di la floro, formacita ek un o plura}
\l{\cel{3}petali, qua envelopas la pistilo e la stamini.- DEFIS.}
\l{}
\l{\sou{korolario}. Konsequo qua esas sequontade la konkluzo di demon-}
\l{\cel{3}stro. - DEFIRS.}
\l{}
\l{\sou{koronala}. (\sou{anat.})\cel{2}Dicesas pri la arterii e veini di la kor-}
\l{\cel{3}dio e di la stomako, dispozita cirkle. -}
\l{}
\l{\sou{koronero.}\cel{2}En Britania : oficiero judiciala di qua la tasko}
\l{\cel{3}konsistas ye inquestar, helpate de jurio, pri la kazi di}
\l{\cel{3}morto violentala, naufrajo, e la deskovri di trezori. -}
\l{}
\l{}
\l{korpo. I. Che la homi e la bestii : ensemblo ek la parti ma-}
\l{\cel{3}teriala qui kompozas organismo en qua rezidas la vivo}
\l{\cel{3}animalala. - II. Omna agregajo ek molekuli materiala.}
\l{\cel{3}- III. (l) Parto precipua di kozo ; (2) ensemblo formaci-}
\l{\cel{3}ta ek kolektajo de personi, de kozi. - DEFIRS.}
\l{}
\l{\sou{korporaciono.}\cel{2}Asociitaro de individui, interligita per re-}
\l{\cel{3}guli, obligesi, privileji komuna. - DEFIRS.}
\l{}
\l{\sou{korporalo}. (\sou{liturgio katol.})\cel{2}Linjo sakrigita quan la sacer-}
\l{\cel{3}doto extensas sur la altaro, lor meso, e sur qua lu pozas}
\l{\cel{3}la kalico e la hostio. - DEFIS.}
\l{}
\l{korpulenta. (\sou{pri homo}). Karakterizata da ampleso, plu o min}
\l{\cel{3}grandega, di sua korpo. - Di qua la tota korpo (anke la}
\l{\cel{3}membri) esas grosa, dika ; muskuloza forte, do robusta e}
\l{\cel{3}sana. - DEFIS.}
\l{}
\l{korpuskulo.\cel{2}Materio-grano infinite mikra. - DEFIS.}
\l{}
\l{\sou{korso.}\cel{2}Aleo uzata kom promeneyo publika, en urbo. DFI.}
\l{}
\l{\sou{korsaj}o. Parto di la robo, qua kovras la busto (di muliero),}
\l{\cel{3}la parto supra, de la shultri til la tayo (do, super la}
\l{\cel{3}jupo). - EF.}
}
